\documentclass[12pt, a4paper]{article}

\usepackage[utf8]{inputenc}
\usepackage[T2A]{fontenc}
\usepackage[russian]{babel}
\usepackage{amsmath, amssymb, amsthm}
\usepackage{geometry}
\usepackage{bm}

\geometry{left=2cm, right=2cm, top=2cm, bottom=2cm}

% Настройки теорем и определений
\theoremstyle{definition}
\newtheorem{definition}{Определение}
\newtheorem{algorithm}{Алгоритм}

\theoremstyle{plain}
\newtheorem{theorem}{Теорема}

\title{Билет №91 \\ \small Определение многомерных КИХ систем. Основные подходы к синтезу. Синтез многомерных КИХ систем с использованием окон. (СпБПУ)}
\author{Сериков Владислав Александрович}
\date{16 мая 2026}

\begin{document}

\maketitle
\tableofcontents
\newpage

\section{Введение и базовые определения}

Многомерная цифровая обработка сигналов имеет дело с сигналами, зависящими от нескольких дискретных независимых переменных (пространственных координат, времени и т.д.). Для общности будем рассматривать $M$-мерное пространство индексов $\mathbf{n} = [n_1, n_2, \dots, n_M]^T \in \mathbb{Z}^M$.

\begin{definition}[Многомерная линейная стационарная система]
Система называется \textit{линейной стационарной} (в пространстве) системой (ЛСС), если ее реакция на произвольный входной сигнал $x[\mathbf{n}]$ полностью описывается многомерной линейной дискретной сверткой с импульсной характеристикой (ИХ) системы $h[\mathbf{n}]$:
\begin{equation}
    y[\mathbf{n}] = x[\mathbf{n}] * h[\mathbf{n}] = \sum_{\mathbf{k} \in \mathbb{Z}^M} h[\mathbf{k}] x[\mathbf{n} - \mathbf{k}].
\end{equation}
\end{definition}

\begin{definition}[Многомерная КИХ-система]
Многомерная система называется системой с \textit{конечной импульсной характеристикой} (КИХ-системой), если ее импульсная характеристика $h[\mathbf{n}]$ имеет конечный носитель. То есть существует такое ограниченное подмножество $\mathcal{S} \subset \mathbb{Z}^M$, что $h[\mathbf{n}] = 0$ для всех $\mathbf{n} \notin \mathcal{S}$.
\end{definition}

Многомерные КИХ-системы обладают рядом важнейших преимуществ: они могут иметь строго линейную фазу (что предотвращает фазовые искажения, критичные для обработки изображений) и гарантированно устойчивы.

\begin{theorem}[Устойчивость многомерных КИХ-систем]
Любая многомерная ЛСС с конечной импульсной характеристикой является абсолютно устойчивой в смысле «ограниченный вход — ограниченный выход» (BIBO).
\end{theorem}
\begin{proof}
Система устойчива по критерию BIBO, если для любого ограниченного входного сигнала ($|x[\mathbf{n}]| \leq B_x < \infty$) выходной сигнал также ограничен ($|y[\mathbf{n}]| \leq B_y < \infty$).
Запишем выражение для модуля выходного сигнала, используя неравенство треугольника:
\begin{equation}
    |y[\mathbf{n}]| = \left| \sum_{\mathbf{k} \in \mathcal{S}} h[\mathbf{k}] x[\mathbf{n} - \mathbf{k}] \right| \leq \sum_{\mathbf{k} \in \mathcal{S}} |h[\mathbf{k}]| \cdot |x[\mathbf{n} - \mathbf{k}]|.
\end{equation}
Поскольку входной сигнал ограничен, вынесем $B_x$ за знак суммы:
\begin{equation}
    |y[\mathbf{n}]| \leq B_x \sum_{\mathbf{k} \in \mathcal{S}} |h[\mathbf{k}]|.
\end{equation}
Так как система является КИХ-системой, множество $\mathcal{S}$ содержит конечное число $N$ элементов. Предполагая, что каждое значение импульсной характеристики конечно (пусть $\max_{\mathbf{k}} |h[\mathbf{k}]| = M_h < \infty$), получаем:
\begin{equation}
    \sum_{\mathbf{k} \in \mathcal{S}} |h[\mathbf{k}]| \leq N \cdot M_h = C < \infty.
\end{equation}
Следовательно, $|y[\mathbf{n}]| \leq B_x \cdot C = B_y < \infty$, что и требовалось доказать.
\end{proof}

\section{Основные подходы к синтезу многомерных КИХ-систем}

Задача синтеза фильтра сводится к нахождению таких коэффициентов $h[\mathbf{n}]$, чтобы частотная характеристика системы $H(e^{j\bm{\omega}})$ наилучшим образом аппроксимировала заданную (идеальную) частотную характеристику $H_d(e^{j\bm{\omega}})$.

К основным подходам к синтезу многомерных (в первую очередь, двумерных) КИХ-фильтров относятся:

\begin{enumerate}
    \item \textbf{Метод окон (Windowing method).} Наиболее прямой метод, основанный на усечении идеальной бесконечной импульсной характеристики с помощью пространственного окна.
    \item \textbf{Метод частотной выборки (Frequency sampling).} Задаются значения частотной характеристики в дискретных точках частотной плоскости, а ИХ вычисляется с помощью обратного дискретного преобразования Фурье.
    \item \textbf{Метод преобразований (Transformation method, преобразование МакКлеллана).} Позволяет преобразовать одномерный оптимальный фильтр в двумерный путем замены переменных (например, $\cos(\omega) \to F(\omega_1, \omega_2)$). Сохраняет свойства оптимальности и линейную фазу.
    \item \textbf{Оптимизационные (минимаксные) методы.} Минимизируют максимальную ошибку аппроксимации в заданной полосе (по критерию Чебышева). В 2D-случае алгоритм Ремеза не имеет прямого аналитического расширения (из-за отсутствия теоремы об альтернансе Чебышева в 2D), поэтому используются алгоритмы линейного программирования или итеративные методы замены.
\end{enumerate}

\section{Синтез многомерных КИХ-систем с использованием окон}

Метод окон заключается в перемножении идеальной бесконечной импульсной характеристики $h_d[\mathbf{n}]$ на функцию окна $w[\mathbf{n}]$, имеющую конечный носитель.

\subsection{Теоретическая база}

\begin{theorem}[Свертка спектров при оконном преобразовании]
Пусть $h[\mathbf{n}] = h_d[\mathbf{n}] \cdot w[\mathbf{n}]$, где $\mathbf{n} \in \mathbb{Z}^M$. Тогда частотная характеристика фильтра $H(\bm{\omega})$ (где $\bm{\omega} \in [-\pi, \pi]^M$) связана с идеальной частотной характеристикой $H_d(\bm{\omega})$ и спектром окна $W(\bm{\omega})$ многомерной круговой сверткой:
\begin{equation}
    H(\bm{\omega}) = \frac{1}{(2\pi)^M} \int_{[-\pi, \pi]^M} H_d(\bm{\theta}) W(\bm{\omega} - \bm{\theta}) \, d\bm{\theta}.
\end{equation}
\end{theorem}
\begin{proof}
Запишем определение многомерного дискретного преобразования Фурье (ДВМПФ) для $h[\mathbf{n}]$:
\begin{equation}
    H(\bm{\omega}) = \sum_{\mathbf{n} \in \mathbb{Z}^M} h[\mathbf{n}] e^{-j \bm{\omega}^T \mathbf{n}} = \sum_{\mathbf{n} \in \mathbb{Z}^M} \left( h_d[\mathbf{n}] w[\mathbf{n}] \right) e^{-j \bm{\omega}^T \mathbf{n}}.
\end{equation}
Выразим $w[\mathbf{n}]$ через обратное ДВМПФ его спектра $W(\bm{\theta})$:
\begin{equation}
    w[\mathbf{n}] = \frac{1}{(2\pi)^M} \int_{[-\pi, \pi]^M} W(\bm{\theta}) e^{j \bm{\theta}^T \mathbf{n}} \, d\bm{\theta}.
\end{equation}
Подставим это выражение в исходную сумму:
\begin{equation}
    H(\bm{\omega}) = \sum_{\mathbf{n}} h_d[\mathbf{n}] \left( \frac{1}{(2\pi)^M} \int_{[-\pi, \pi]^M} W(\bm{\theta}) e^{j \bm{\theta}^T \mathbf{n}} \, d\bm{\theta} \right) e^{-j \bm{\omega}^T \mathbf{n}}.
\end{equation}
Поскольку сигналы абсолютно суммируемы (окно финитно), мы можем поменять местами интегрирование и суммирование (теорема Фубини):
\begin{equation}
    H(\bm{\omega}) = \frac{1}{(2\pi)^M} \int_{[-\pi, \pi]^M} W(\bm{\theta}) \left( \sum_{\mathbf{n}} h_d[\mathbf{n}] e^{-j (\bm{\omega} - \bm{\theta})^T \mathbf{n}} \right) d\bm{\theta}.
\end{equation}
Заметим, что выражение в скобках является определением ДВМПФ для $h_d[\mathbf{n}]$, сдвинутого на частоту $\bm{\theta}$:
\begin{equation}
    \sum_{\mathbf{n}} h_d[\mathbf{n}] e^{-j (\bm{\omega} - \bm{\theta})^T \mathbf{n}} = H_d(\bm{\omega} - \bm{\theta}).
\end{equation}
Отсюда получаем итоговое выражение:
\begin{equation}
    H(\bm{\omega}) = \frac{1}{(2\pi)^M} \int_{[-\pi, \pi]^M} W(\bm{\theta}) H_d(\bm{\omega} - \bm{\theta}) \, d\bm{\theta}.
\end{equation}
\end{proof}

Данная теорема объясняет \textit{эффект Гиббса}: из-за того, что спектр реального окна $W(\bm{\omega})$ имеет главный лепесток конечной ширины и боковые лепестки, в результирующей частотной характеристике $H(\bm{\omega})$ появляются пульсации в полосах пропускания и подавления, а также расширяется полоса перехода.

\subsection{Алгоритм синтеза}

\begin{algorithm}[Синтез двумерного КИХ-фильтра методом окон] \leavevmode
\begin{description}
    \item[Шаг 1. Формирование идеальной частотной характеристики.] \hfill \\
    Задается функция $H_d(\omega_1, \omega_2)$, которая равна $1$ в полосе пропускания и $0$ в полосе подавления.

    \item[Шаг 2. Вычисление идеальной ИХ.] \hfill \\
    Аналитически или численно вычисляется обратное преобразование Фурье:
    \begin{equation}
        h_d[n_1, n_2] = \frac{1}{4\pi^2} \int_{-\pi}^{\pi} \int_{-\pi}^{\pi} H_d(\omega_1, \omega_2) e^{j(\omega_1 n_1 + \omega_2 n_2)} \, d\omega_1 d\omega_2.
    \end{equation}

    \item[Шаг 3. Выбор функции окна.] \hfill \\
    Исходя из требуемого подавления в полосе задерживания и допустимой ширины переходной полосы, выбирается форма 2D-окна. Окно может быть:
    \begin{itemize}
        \item Сепарабельным (прямоугольным): $w[n_1, n_2] = w_{1D}[n_1] \cdot w_{1D}[n_2]$;
        \item Круговосимметричным: $w[n_1, n_2] = w_{1D}(\sqrt{n_1^2 + n_2^2})$.
    \end{itemize}

    \item[Шаг 4. Усечение и применение окна.] \hfill \\
    Формируется конечная ИХ фильтра:
    \begin{equation}
        h[n_1, n_2] = h_d[n_1, n_2] \cdot w[n_1, n_2].
    \end{equation}

    \item[Шаг 5. Сдвиг (опционально).] \hfill \\
    Для реализации каузальной системы носитель смещается в область неотрицательных индексов (если фильтр реализуется во временной области). При обработке пространственных данных (изображений) сдвиг обычно не требуется, так как система может быть некаузальной.
\end{description}
\end{algorithm}

\subsection{Минимальный пример: Идеальный круговой ФНЧ}

Рассмотрим применение алгоритма для синтеза двумерного фильтра нижних частот (ФНЧ) с круговой симметрией и частотой среза $\omega_c$.

\textbf{Шаг 1.} Идеальная АЧХ задается как цилиндр:
\begin{equation}
    H_d(\omega_1, \omega_2) =
    \begin{cases}
      1, & \text{если } \sqrt{\omega_1^2 + \omega_2^2} \leq \omega_c, \\
      0, & \text{иначе}.
   \end{cases}
\end{equation}

\textbf{Шаг 2.} Вычисляем идеальную ИХ через обратное ДВМПФ. Переходя к полярным координатам в частотной ($\rho, \phi$) и пространственной ($r, \theta$) областях (где $r = \sqrt{n_1^2 + n_2^2}$), интеграл сводится к:
\begin{equation}
    h_d[n_1, n_2] = \frac{1}{2\pi} \int_0^{\omega_c} \rho J_0(\rho r) \, d\rho = \frac{\omega_c}{2\pi \sqrt{n_1^2 + n_2^2}} J_1 \left(\omega_c \sqrt{n_1^2 + n_2^2}\right),
\end{equation}
где $J_0$ и $J_1$ — функции Бесселя первого рода нулевого и первого порядка соответственно. В точке начала координат по правилу Лопиталя: $h_d[0, 0] = \frac{\omega_c^2}{4\pi}$.

\textbf{Шаг 3 и 4.} Выбираем двумерное окно Хэмминга $w_H[n_1, n_2]$ размером $(2N+1) \times (2N+1)$ и умножаем на него $h_d[n_1, n_2]$. Полученная матрица коэффициентов $h[n_1, n_2]$ будет служить маской (ядром свертки) фильтра для цифрового изображения.

\end{document}
